% =============================================================================
%  lec06.tex —— 第 6 讲 极限与余极限：定义（底本 PDF 第 50–57 页）
% =============================================================================
\lecture{6}{极限与余极限：定义}
\begin{center}\small\color{black!55}（底本 PDF 第 50–57 页）\end{center}

本讲我们定义 $\infty$-(余)极限，即 $\infty$-范畴中的（余）极限。

\section{定义的想法}

\nitem\label{L6-1.1} 在经典范畴论中，我们有如下极限的刻画：
\begin{center}
\emph{一个图 $\bar u\colon K^{\triangleleft}\to\cC$ 是极限图，当且仅当它对应于切片范畴 $\cC_{/u}$ 中的一个终对象，其中 $u\colon K\to\cC$。}
\end{center}
只要完成下面两件事，我们就将用完全相同的方式定义 $\infty$-极限：
\begin{enumerate}[label=(\arabic*),leftmargin=2.4em,itemsep=0.1em]
\item 在 $\infty$-范畴中定义终对象与始对象；
\item 定义切片与余切片 $\infty$-范畴。
\end{enumerate}

\section{终对象与始对象}

\nitem\label{L6-2.1} 回忆：$\infty$-群胚是\textbf{可缩的}（contractible），是指它等价于 $[0]$。

\begin{definition}\label{L6-2.2}
设 $\cC$ 是 $\infty$-范畴，$x\in\cC$ 是对象。
\begin{itemize}[leftmargin=2.2em,itemsep=0.1em]
\item 若对任意对象 $y\in\cC$，$\infty$-群胚 $\operatorname{Maps}_{\cC}(y,x)$ 可缩，则称 $x\in\cC$ 为\textbf{终对象}（final object）；
\item 若对任意对象 $y\in\cC$，$\infty$-群胚 $\operatorname{Maps}_{\cC}(x,y)$ 可缩，则称 $x\in\cC$ 为\textbf{始对象}（initial object）。
\end{itemize}
\end{definition}

\yizhu{原文用 final object，另一些文献用 terminal object，二者同义；本书一律作 final object，中译统一为“终对象”。相应地，initial object 统一为“始对象”。\textbf{[来源：B]}}

\nitem\label{L6-2.3} 注意 $x\in\cC$ 是始对象，当且仅当它在 $\cC^{\op}$ 中对应的对象是终对象。因此下面我们集中讨论终对象。

\begin{exercise}\label{L6-2.4}
是终对象这一性质在等价下不变：
\begin{enumerate}[label=(\arabic*),leftmargin=2.4em,itemsep=0.1em]
\item 若 $x\to y$ 是 $\cC$ 中的同构，则 $x$ 是终对象当且仅当 $y$ 是终对象；
\item 若 $F\colon\cC\to\cC'$ 是 $\infty$-范畴之间的等价，则 $x\in\cC$ 是终对象当且仅当 $F(x)\in\cC'$ 是终对象。
\end{enumerate}
\end{exercise}

\begin{exercise}\label{L6-2.5}
$C^{\triangleright}$ 的顶点是终对象。
\end{exercise}

\begin{example}\label{L6-2.6}
空范畴在 $\Cat_\infty$ 中是始对象；单点 $[0]$ 在 $\Cat_\infty$ 中是终对象。
\end{example}

\begin{proposition}\label{L6-2.7}
$\infty$-范畴 $\cC$ 中由全体终对象构成的全子 $\infty$-范畴，或者为空，或者是可缩的 $\infty$-群胚。换言之，$\infty$-范畴 $\cC$ 中的终对象若存在，则本质唯一。
\end{proposition}

\begin{proof}
我们可以把 $\cC$ 换成由全体终对象构成的全子 $\infty$-范畴，并假设 $\cC$ 中每个对象都是终对象。我们要证明 $\cC$ 或者为空，或者是可缩的 $\infty$-群胚。事实上，若 $\cC$ 非空，则函子 $\cC\to[0]$ 全忠实且本质满。
\end{proof}

\begin{exercise}\label{L6-2.8}
设 $\cC$ 是 $\infty$-范畴，$x$ 是对象。若 $x$ 是终对象，则它在 $\cC\to\mathrm{h}\cC$ 下的像也是终对象。若 $\cC$ 有终对象，则逆命题成立。
\end{exercise}

\begin{warning}\label{L6-2.9}
一般来说，函子 $\cC\to\mathrm{h}\cC$ 无法探测终对象。例如，下述命题表明：不可缩的 $\infty$-群胚 $\cC$ 没有终对象，但当 $\cC$ 单连通时 $\mathrm{h}\cC\simeq[0]$。
\end{warning}

\begin{proposition}\label{L6-2.10}
$\infty$-群胚 $\cC$ 含有终对象 $x$，当且仅当 $\cC$ 可缩。
\end{proposition}

\begin{proof}[证明梗概]
\textbf{梗概 1.}\quad 对带基点的空间 $X$，$\pi_n(\Omega X)\simeq\pi_{n+1}(X)$。

\textbf{梗概 2.}\quad 函子 $[0]\xrightarrow{x}\cC$ 全忠实且本质满。
\end{proof}

\nitem\label{L6-2.11} 事实上，我们有下述结果，它将在 §4 中证明。

\begin{proposition}\label{L6-2.12}
若 $\infty$-范畴 $\cC$ 含有终对象 $x$，则它是弱可缩的\footnote{换言之，把 $\cC$ 中所有态射形式地可逆化之后，$\cC$ 成为可缩的 $\infty$-群胚。见第 5 讲 §3.7--3.10。}。
\end{proposition}

\section{切片与余切片：一般情形}

\nitem\label{L6-3.1} 在第 4 讲 §8 中，对拟范畴 $\cC$ 与对象 $x\in\cC$，我们定义了切片与余切片拟范畴 $\cC_{/x}$ 与 $\cC_{x/}$。本节把这一构造推广到 $\infty$-范畴 $\cC$ 中的一般图 $u\colon K\to\cC$，以定义 $\infty$-范畴 $\cC_{/u}$ 与 $\cC_{u/}$。

\nitem\label{L6-3.2} 我们先通过拟范畴来定义它们。

\begin{definition}\label{L6-3.3}
设 $u\colon K\to X$ 是 $\Set_\Delta$ 中的态射。$X$ 在 $u$ 上的\textbf{切片单纯集}是如下定义的单纯集 $X_{/u}$\footnote{联 $Y\star K$ 的定义见第 4 讲定义 8.5。}：
\[
\Hom(Y,X_{/u})\simeq\Hom(Y\star K,X)\underset{\Hom(K,X)}{\times}\{u\}.
\]
对偶地，$X$ 在 $u$ 下的\textbf{余切片单纯集}是如下定义的单纯集 $X_{u/}$：
\[
\Hom(Y,X_{u/})\simeq\Hom(K\star Y,X)\underset{\Hom(K,X)}{\times}\{u\}.
\]
\end{definition}

\nitem\label{L6-3.4} 注意 $(X_{u/})^{\op}\simeq (X^{\op})_{/u^{\op}}$。因此下面我们集中讨论切片构造。

\begin{proposition}\label{L6-3.5}
（Ker.018F）设 $u\colon K\to\cC$ 是 $\Set_\Delta$ 中的态射，且 $\cC$ 是拟范畴。则 $\cC_{/u}$ 是拟范畴。
\end{proposition}

\begin{proposition}\label{L6-3.6}
（Ker.02GL, 02NC, 02NR）切片构造在等价下不变：
\begin{enumerate}[label=(\arabic*),leftmargin=2.4em,itemsep=0.1em]
\item 若 $F\colon\cC\to\cC'$ 是拟范畴之间的等价，则对任意图 $u\colon K\to\cC$，诱导函子 $\cC_{/u}\to\cC'_{/F\circ u}$ 是等价；
\item 若 $v\colon K'\to K$ 是单纯集之间的范畴等价（categorical equivalence），则对拟范畴 $\cC$ 中的任意图 $u\colon K\to\cC$，限制函子 $\cC_{/u}\to\cC_{/u\circ v}$ 是等价。
\end{enumerate}
\end{proposition}

\yizhu{categorical equivalence 译“范畴等价”，指 $\Set_\Delta$ 中 Joyal 模型结构下的弱等价，即（在同伦意义下）可逆到某个范畴的态射；它与“$\infty$-范畴的等价”（equivalence of $\infty$-categories，即函子之间的等价）不在同一层次上。\textbf{[来源：B]}}

\nitem\label{L6-3.7} 由第 5 讲 §9 的哲学，对 $\infty$-范畴 $\cC$ 中的图 $u\colon K\to\cC$，存在良定义的 $\infty$-范畴 $\cC_{/u}$ 与 $\cC_{u/}$；并且当 $K$ 换成 $\infty$-范畴时同样成立。

我们称它们为 $\cC$ 关于 $u$ 的\textbf{切片/余切片 $\infty$-范畴}。

\begin{exercise}\label{L6-3.8}
设 $u\colon K\to\cC$ 是 $\infty$-范畴 $\cC$ 中的图，$v\colon J\to\cC_{/u}$ 是 $\cC_{/u}$ 中的图。证明 $(\cC_{/u})_{/v}\simeq\cC_{/w}$，其中 $w\colon J\star K\to\cC$ 是对应于 $v$ 的图。
\end{exercise}

\begin{construction}\label{L6-3.9}
设 $u\colon K\to X$ 是 $\Set_\Delta$ 中的态射。由 Yoneda 引理，态射 $Y\to Y\star K$ 诱导出态射 $X_{/u}\to X$，称为从 $X_{/u}$ 到 $X$ 的\textbf{遗忘态射}。

设 $u\colon K\to\cC$ 是 $\infty$-范畴 $\cC$ 中的图。上述构造给出遗忘函子 $\cC_{/u}\to\cC$。
\end{construction}

\nitem\label{L6-3.10} 我们将在 §5 中证明下述结果。

\begin{proposition}\label{L6-3.11}
对 $\infty$-范畴 $\cC$ 中的任意图 $u\colon K\to\cC$，遗忘函子 $\cC_{/u}\to\cC$ 是\textbf{保守的}（conservative）\footnote{这是指：$\cC_{/u}$ 中的态射是同构，当且仅当它在 $\cC$ 中的像是同构。}。
\end{proposition}

\begin{proposition}\label{L6-3.12}
（Ker.018F, Ker.00TE）设 $u\colon K\to\cC$ 是 $\infty$-范畴 $\cC$ 中的图。对 $\cC$ 中的态射 $f\colon x\to y$ 以及 $y$ 的提升 $\tilde y\in\cC_{/u}$，存在本质唯一的态射 $\tilde x\to\tilde y$ 提升 $f$。
\end{proposition}

\begin{remark}\label{L6-3.13}
该命题的确切含义是：
\begin{enumerate}[label=(\arabic*),leftmargin=2.4em,itemsep=0.1em]
\item 函子
\[
\Fun([1],\cC_{/u})\to\Fun(\{1\},\cC_{/u})\underset{\Fun(\{1\},\cC)}{\times}\Fun([1],\cC)
\]
是 $\infty$-范畴之间的等价。

这里纤维积是在 $\infty$-范畴 $\Cat_\infty$ 中计算的，后者将在下一节定义。

或者，我们可以把 $\cC$ 实现为拟范畴 $\cC$，此时该命题说的是\footnote{为这些讲义的目的，我们尽可能把 Lurie 书中的结果翻译成不依赖模型的表述。因此，本讲义中陈述的内容与所引文献中的陈述并不完全相同；要比较它们，人们往往需要关于模型范畴 $\Set^{\mathrm{Joyal}}_{\Delta}$ 的额外知识。例如，要从（2）推出（1），需要知道（1）中的纤维积可以实现为（2）中的纤维积。由 HTT.4.2.4.1（将在以后的讲中讨论），前者可以计算为单纯模型范畴 $\Set^{\mathrm{Joyal}}_{\Delta}$ 中相应的同伦纤维积；由 HTT.A.2.4.4（的对偶），该同伦纤维积与朴素纤维积一致，因为 $\Fun(\Delta^1,\cC)\to\Fun(\{1\},\cC)$ 是范畴纤维化（HTT.2.4.6.5），且三个项在 $\Set^{\mathrm{Joyal}}_{\Delta}$ 中都是纤维性的。}\label{fn:L6-4}：
\item 函子
\[
\Fun(\Delta^1,\cC_{/u})\to\Fun(\{1\},\cC_{/u})\underset{\Fun(\{1\},\cC)}{\times}\Fun(\Delta^1,\cC)
\]
是无环 Kan 纤维化（acyclic Kan fibration）。

这里纤维积是在通常范畴 $\mathrm{QCat}$ 中计算的。
\end{enumerate}
满足命题~\ref{L6-3.12} 中提升性质的函子 $\cD\to\cC$（例如 $\cC_{/u}\to\cC$）称为\textbf{右纤维化}（right fibration）\footnote{更确切地说，这类函子由拟范畴之间的右纤维化建模。}。对偶地，$\cC_{u/}\to\cC$ 是\textbf{左纤维化}（left fibration）。我们将在以后的讲中系统地研究这些纤维化。
\end{remark}

\nitem\label{L6-3.14} 下面讨论切片构造关于 $u$ 的函子性。

\nitem\label{L6-3.15}\textbf{命题—构造}\quad 设 $K$ 是单纯集，$\cC$ 是 $\infty$-范畴。设 $u_0\to u_1$ 是 $\Fun(K,\cC)$ 中的态射，$u\colon K\times\Delta^1\to\cC$ 是相应的图。考虑函子
\[
\cC_{/u_0}\leftarrow\cC_{/u}\to\cC_{/u_1},
\]
它们由沿 $K\rightrightarrows K\times\Delta^1$ 的限制得到。
\begin{enumerate}[label=(\arabic*),leftmargin=2.4em,itemsep=0.1em]
\item 函子 $\cC_{/u_0}\leftarrow\cC_{/u}$ 是等价。
\item 若 $u_0\to u_1$ 是 $\Fun(K,\cC)$ 中的同构，则函子 $\cC_{/u}\to\cC_{/u_1}$ 是等价。
\end{enumerate}
将等价 $\cC_{/u_0}\leftarrow\cC_{/u}$ 可逆化，我们得到与到 $\cC$ 的遗忘函子相容的典范函子 $\cC_{/u_0}\to\cC_{/u_1}$。

\begin{proof}
（1）留作下一讲的练习。为证明（2），考虑拟范畴 $\cI\coloneqq\mathrm{N}_{\bullet}(I)$，其中 $I$ 是满足下述条件的通常范畴：
\begin{itemize}[leftmargin=2.2em,itemsep=0.1em]
\item $I$ 中有两个对象；
\item 对 $x,y\in\Ob(I)$，存在唯一的态射 $x\to y$。
\end{itemize}
由于 $u_0\to u_1$ 是同构，我们可以找到态射 $v$ 使下图交换：
\begin{center}
\begin{tikzcd}[column sep=3.6em,row sep=2.8em]
K\times\Delta^1 \arrow[r,"u"] \arrow[d, hook, "\subset"'] & \cC \\
K\times\cI \arrow[ur, dashed, "v"] &
\end{tikzcd}
\end{center}
于是我们有下述交换图
\begin{center}
\begin{tikzcd}[column sep=3.4em,row sep=3.0em]
& \cC_{/v} \arrow[dl,"{p'}"'] \arrow[d,"r"] \arrow[dr,"{q'}"] & \\
\cC_{/u_0} & \cC_{/u} \arrow[l,"p"'] \arrow[r,"q"] & \cC_{/u_1}
\end{tikzcd}
\end{center}
由（1），$p$ 是等价。由命题~\ref{L6-3.6}(2)，$p'$ 与 $q'$ 都是等价。由此 $r$、从而 $q$ 是等价。
\end{proof}

\begin{exercise}\label{L6-3.16}
设 $K$ 是单纯集，$\cC$ 是拟范畴。你能构造一个函子 $\Fun(K,\cC)\to\mathrm{QCat}$，$u\mapsto\cC_{/u}$ 吗？这里 $\mathrm{QCat}$ 是拟范畴的通常范畴。
\end{exercise}

\section{终对象与切片构造}

\begin{proposition}\label{L6-4.1}
（Ker.02J2）设 $\cC$ 是 $\infty$-范畴，$f\colon x\to y$ 是态射。则下述条件等价：
\begin{itemize}[leftmargin=2.2em,itemsep=0.1em]
\item $f$ 是同构；
\item 由 $f$ 给出的 $\cC_{x/}$ 中的对象是始对象；
\item 由 $f$ 给出的 $\cC_{/y}$ 中的对象是终对象。
\end{itemize}
\end{proposition}

\begin{proposition}\label{L6-4.2}
设 $\cC$ 是 $\infty$-范畴，$x\in\cC$ 是对象。则下述条件等价：
\begin{enumerate}[label=(\arabic*),leftmargin=2.4em,itemsep=0.1em]
\item 对象 $x$ 是终对象；
\item 函子 $\cC_{/x}\to\cC$ 是等价；
\item 函子 $\cC_{/x}\to\cC$ 是平凡 Kan 纤维化（当把 $\cC$ 实现为拟范畴时）。
\end{enumerate}
\end{proposition}

\begin{proof}
（1）$\Leftarrow$（2）$\Leftarrow$（3）显然。（1）$\Rightarrow$（3）即 Ker.02HF。
\end{proof}

\begin{proof}[命题~\ref{L6-2.12} 的证明]
按定义，我们只需证明：对任意 $\infty$-群胚 $\cD$，函子
\[
-\circ\pi\colon\Fun([0],\cD)\to\Fun(\cC,\cD)
\]
是等价，其中 $\pi\colon\cC\to[0]$ 是投影函子。注意该函子有显然的左逆
\[
-\circ x\colon\Fun(\cC,\cD)\to\Fun([0],\cD),
\]
它由函子 $x\colon[0]\to\cC$ 诱导。于是只需构造从
\begin{equation}\label{L6-eq-4.1}
\Fun(\cC,\cD)\xrightarrow{\id}\Fun(\cC,\cD)
\end{equation}
到
\begin{equation}\label{L6-eq-4.2}
\Fun(\cC,\cD)\xrightarrow{-\circ x}\Fun([0],\cD)\xrightarrow{-\circ\pi}\Fun(\cC,\cD)
\end{equation}
的可逆自然变换。

我们断言：任何这样的自然变换都自动可逆。为证明该断言，注意 $\Fun(\cC,\cD)$ 是 $\infty$-群胚，因为 $\cD$ 是 $\infty$-群胚。因此
\[
\Fun(\Fun(\cC,\cD),\Fun(\cC,\cD))
\]
也是 $\infty$-群胚，这蕴含该断言。

剩下要构造从~\eqref{L6-eq-4.1} 到~\eqref{L6-eq-4.2} 的自然变换。为此，只需构造从 $\id_{\cC}$ 到复合 $\cC\xrightarrow{\pi}[0]\xrightarrow{x}\cC$ 的自然变换。考虑显然的函子
\[
\cC\leftarrow\cC_{/x}\to\Fun([1],\cC).
\]
由命题~\ref{L6-4.2}，向左的函子是等价。于是可以将其可逆化，得到函子 $\cC\to\Fun([1],\cC)$，它对应于函子 $[1]\to\Fun(\cC,\cC)$。展开定义，这给出从 $\id_{\cC}$ 到 $\cC\xrightarrow{\pi}[0]\xrightarrow{x}\cC$ 的自然变换。
\end{proof}

\section{钝联}

\nitem\label{L6-5.1} 为证明命题~\ref{L6-3.11}，我们需要切片 $\infty$-范畴的下述等价构造。设 $u\colon K\to\cC$ 是拟范畴 $\cC$ 中的图。如同第 4 讲 §9.1，我们有典范函子
\begin{equation}\label{L6-eq-5.1}
\cC_{/u}\to\cC\underset{\Fun(K\times\{0\},\cC)}{\times}\Fun(K\times\Delta^1,\cC)\underset{\Fun(K\times\{1\},\cC)}{\times}\{u\}.
\end{equation}
在对象层面，该函子把图 $\bar u\colon K^{\triangleleft}\to\cC$ 送到复合 $K\times\Delta^1\to K^{\triangleleft}\to\cC$，它对应于自然变换 $\underline{x}\to u$，其中 $\underline{x}\colon K\to\cC$ 是取值 $x\coloneqq\bar u(\ast)$ 的常函子。

为简化记号，我们作如下定义。

\begin{definition}\label{L6-5.2}
设 $J$ 与 $K$ 是单纯集。定义 $J$ 与 $K$ 的\textbf{钝联}（blunt join）为
\[
J\Diamond K\coloneqq J\underset{J\times K\times\{0\}}{\sqcup}(J\times K\times\Delta^1)\underset{J\times K\times\{1\}}{\sqcup}K.
\]
\end{definition}

\yizhu{“钝联”译自 blunt join，原文未加解释。就其定义而言，它把 $J$ 与 $K$ 之间原本由单形联 $J\star K$ 逐层粘合出的部分，整体替换为 $J\times K\times\Delta^1$，故名“钝”；记号 $\Diamond$ 也特意区别于通常的联 $\star$。\textbf{[来源：B]}}

\begin{warning}\label{L6-5.3}
拟范畴的钝联未必是拟范畴。
\end{warning}

\begin{exercise}\label{L6-5.4}
证明 $\Delta^0\Diamond\Delta^1$ 不是拟范畴。
\end{exercise}

\nitem\label{L6-5.5} 利用上述记号，函子~\eqref{L6-eq-5.1} 可以改写为
\begin{equation}\label{L6-eq-5.2}
\cC_{/u}\to\Fun(\Delta^0\Diamond K,\cC)\underset{\Fun(K,\cC)}{\times}\{u\}=:\cC\underset{\Fun(K,\cC)}{\overset{\to}{\times}}\{u\}=:\cC^{/u}.
\end{equation}
对偶地，我们有典范函子
\begin{equation}\label{L6-eq-5.3}
\cC_{u/}\to\{u\}\underset{\Fun(K,\cC)}{\times}\Fun(K\Diamond\Delta^0,\cC)=:\{u\}\underset{\Fun(K,\cC)}{\overset{\to}{\times}}\cC=:\cC^{u/}.
\end{equation}

\begin{exercise}\label{L6-5.6}
存在唯一的态射 $J\Diamond K\to J\star K$ 使下图交换：
\begin{center}
\begin{tikzcd}[column sep=4.2em,row sep=2.8em]
& J\Diamond K \arrow[dr, dashed] & \\
J\sqcup K \arrow[ur] \arrow[rr] & & J\star K
\end{tikzcd}
\end{center}
\end{exercise}

\begin{proposition}\label{L6-5.7}
（Ker.01HU, 01HV, 01HW）函子 $-\star-$ 与 $-\Diamond-$ 在两边都保持范畴等价，且态射 $J\Diamond K\to J\star K$ 是范畴等价。
\end{proposition}

\begin{remark}\label{L6-5.8}
因此，我们在 $\infty$-范畴上得到一个良定义的二元运算，它可以由 $-\star-$ 建模。
\end{remark}

\begin{exercise}\label{L6-5.9}
设 $\cC$ 与 $\cD$ 是 $\infty$-范畴。用泛性质刻画 $\cC\star\cD$。
\end{exercise}

\begin{proposition}\label{L6-5.10}
（Ker.01KU）函子~\eqref{L6-eq-5.2} 与~\eqref{L6-eq-5.3} 可逆。
\end{proposition}

\begin{remark}\label{L6-5.11}
可以证明，\eqref{L6-eq-5.2} 中在 $\mathrm{QCat}$ 中取的纤维积，同时也计算 $\Set^{\mathrm{Joyal}}_{\Delta}$ 中的同伦纤维积，从而也计算 $\mathrm{QCat}$ 中的纤维积（见脚注~\ref{fn:L6-4}）。因此，对 $\infty$-范畴中的图 $u\colon K\to\cC$，我们有典范等价
\begin{equation}\label{L6-eq-5.4}
\cC_{/u}\xrightarrow{\simeq}\Fun(\Delta^0\Diamond K,\cC)\underset{\Fun(K,\cC)}{\times}\{u\}\xleftarrow{\simeq}\Fun(\Delta^0\star K,\cC)\underset{\Fun(K,\cC)}{\times}\{u\}.
\end{equation}
\end{remark}

\begin{proof}[命题~\ref{L6-3.11} 的证明]
我们只需对 $\cC^{/u}$ 证明类似的断言。按定义，该拟范畴中的态射是函子 $F_1,F_2\colon K\times\Delta^1\to\cC$ 之间的自然变换 $\alpha\colon F_1\to F_2$，使得
\begin{enumerate}[label=(\arabic*),leftmargin=2.4em,itemsep=0.1em]
\item $\alpha|_{K\times\{0\}}$ 是两个常函子 $\underline{x}_1,\underline{x}_2\colon K\to\cC$ 之间的常自然变换；
\item $\alpha|_{K\times\{1\}}$ 是函子 $u\colon K\to\cC$ 处的恒等自然变换。
\end{enumerate}
注意遗忘函子 $\cC^{/u}\to\cC$ 把 $\alpha$ 送到 (1) 中的态射 $x_1\to x_2$。由第 5 讲定理 2.3，$\alpha$ 可逆当且仅当 $x_1\to x_2$ 可逆。
\end{proof}

\begin{exercise}\label{L6-5.12}
设 $\tilde u\colon J\times K\to\cC$ 是拟范畴 $\cC$ 中的图。将 $\Fun(K,\cC)$ 中相应的图记作 $u\colon J\to\Fun(K,\cC)$。构造典范同构
\begin{equation}\label{L6-eq-5.5}
\cC^{/\tilde u}\xrightarrow{\simeq}\Fun(K,\cC)^{/u}\underset{\Fun(K,\cC)}{\times}\cC.
\end{equation}
\end{exercise}

\section{极限与余极限的定义}

\begin{definition}\label{L6-6.1}
设 $K$ 是单纯集，$\cC$ 是 $\infty$-范畴。若图 $\bar u\colon K^{\triangleleft}\to\cC$ 对应的 $\cC_{/u}$ 中的对象是终对象，则称 $\bar u$ 是\textbf{极限图}（limit diagram），其中 $u\coloneqq\bar u|_K$ 是 $\bar u$ 在 $K$ 上的限制。此时也称 $\bar u$ 把 $\bar u(\ast)$ 实现为图 $u\colon K\to\cC$ 的\textbf{极限}（limit）。

对偶地，若图 $\bar u\colon K^{\triangleright}\to\cC$ 对应的 $\cC_{u/}$ 中的对象是始对象，则称 $\bar u$ 是\textbf{余极限图}（colimit diagram）。此时也称 $\bar u$ 把 $\bar u(\ast)$ 实现为图 $u\colon K\to\cC$ 的\textbf{余极限}（colimit）。
\end{definition}

\nitem\label{L6-6.2} 注意图 $\bar u\colon K^{\triangleright}\to\cC$ 是余极限图，当且仅当 $\bar u^{\op}\colon(K^{\op})^{\triangleleft}\to\cC^{\op}$ 是极限图。因此下面我们集中讨论极限图。

\begin{remark}\label{L6-6.3}
对图 $\bar u\colon K^{\triangleleft}\to\cC$ 而言，是极限图是一个性质。然而，对对象 $x$ 而言，是给定图 $u\colon K\to\cC$ 的极限则是一个结构。
\end{remark}

\begin{exercise}\label{L6-6.4}
空图的极限是终对象。
\end{exercise}

\nitem\label{L6-6.5} 命题~\ref{L6-2.7} 允许我们谈论图 $u\colon K\to\cC$ 的\textbf{极限对象}，只要我们把它延拓后的图 $\bar u$ 作为定义中数据的一部分。极限对象有各种标准记号：
\[
\lim u,\qquad \lim_K u,\qquad \lim_{y\in K} u(y),\ \cdots
\]
使用这些记号时，我们总把它看作 $\cC$ 中的对象，并配以沿遗忘函子 $\cC_{/u}\to\cC$ 的典范提升。

特别地，有典范态射
\[
\mathrm{ev}_y\colon \lim u\to u(y),\quad y\in K,
\]
称为\textbf{求值态射}（evaluating morphisms）。对偶地，我们有\textbf{插入态射}（inserting morphisms）
\[
\mathrm{ins}_y\colon u(y)\to\colim u,\quad y\in K.
\]

\appendixsection{扭箭头}

\begin{exercise}\label{L6-A.1}
设 $\cC$ 是通常的范畴。构造一个范畴 $\mathrm{Tw}(\cC)$，使得
\begin{itemize}[leftmargin=2.2em,itemsep=0.1em]
\item $\mathrm{Tw}(\cC)$ 的对象是 $\cC$ 中的态射 $x\to y$；
\item 指派 $(x\to y)\mapsto x$ 是反变的，即由一个函子 $\mathrm{Tw}(\cC)\to\cC^{\op}$ 给出；
\item 指派 $(x\to y)\mapsto y$ 是协变的，即由一个函子 $\mathrm{Tw}(\cC)\to\cC$ 给出。
\end{itemize}
我们称它为 $\cC$ 中的\textbf{扭箭头范畴}（category of twisted arrows）。
\end{exercise}

\begin{exercise}\label{L6-A.2}
把上述构造推广到 $\infty$-范畴。
\end{exercise}

\begin{exercise}\label{L6-A.3}
构造一个函子 $\mathrm{Tw}(\Cat_\infty)\to\Cat_\infty$，它把扭箭头 $K\xrightarrow{u}\cC$ 送到 $\cC_{/u}$。
\end{exercise}

\nitem\label{L6-A.4} \textbf{建议阅读}\quad Ker.03GB。
